\section{Nonsemiclassical quantum magic square in a corner APP}
\label{sec:cpe-qms-no-semiclasico}

En \(\mathbb C^4\), with base \(e_1,e_2,e_3,e_4\), considerese the arreglo of
vectors
\[
\begin{array}{cccc}
e_1&e_2&e_3&e_4\\
e_2&e_1&e_4&e_3\\
(e_3+e_4)/\sqrt2&(e_3-e_4)/\sqrt2&(e_1+e_2)/\sqrt2&(e_1-e_2)/\sqrt2\\
(e_3-e_4)/\sqrt2&(e_3+e_4)/\sqrt2&(e_1-e_2)/\sqrt2&(e_1+e_2)/\sqrt2.
\end{array}
\]
Let \(P_{ij}\) be the rank-one projector onto the vector located in
cell \((i,j)\). Each row and each column is an orthonormal basis, hence
\begin{equation}
 \sum_jP_{ij}=I_4,
 \qquad \sum_iP_{ij}=I_4.
 \label{eq:cpe-qms4}
\end{equation}
To realize it in the APP fiber of dimension \(9\), fix
\begin{equation}
 A_{ij}:=P_{ij}\oplus\frac14I_5\in M_9(\mathbb C).
 \label{eq:cpe-qms9}
\end{equation}

\begin{theorem}[No certified semiclassicality]
\label{thm:cpe-qms-no-semiclasico}
The matrix \(A=(A_{ij})_{i,j=1}^4\) is a quantum magic square and does not admit
a semiclassical decomposition
\[
 A_{ij}=\sum_{\pi:\,\pi(i)=j}Q_\pi,
 \qquad Q_\pi\succeq0,
 \qquad\sum_\pi Q_\pi=I_9.
\]
\end{theorem}

\begin{proof}
The row and column identities follow from
\eqref{eq:cpe-qms4}; four copies of \(I_5/4\) sum to \(I_5\). If a
semiclassical decomposition existed, its compression to the corner \(\mathbb C^4\)
would be a joint measurement of the PVM defined by the rows
\((P_{ij})_j\). Jointly measurable projective measurements commute.
However,
\[
 \left[
 |e_1\rangle\langle e_1|,
 \left|\frac{e_1+e_2}{\sqrt2}\right\rangle
 \left\langle\frac{e_1+e_2}{\sqrt2}\right|
 \right]
 =\frac12
 \begin{pmatrix}0&1\\-1&0\end{pmatrix}
 \oplus0_2\ne0.
\]
The contradiction proves nonsemiclassicality. The compression simultaneously
provides an analytic witness to the infeasibility of the semidefinite problem.
\end{proof}

Each row defines a UCP channel.
\[
 \Phi_i:\ell^\infty_4\to M_9,
 \qquad \Phi_i(\delta_j)=A_{ij},
\]
and each column defines another. Thus, the magic normalizations translate into
completely positive maps with explicit domain and codomain; a symbol
\(J(\operatorname{id})\) is not used without first fixing those types.

\section{Conical route system and internal genealogical universality}
\label{sec:cpe-s-system-universalidad}

For each finite profile \(c\), let \(G_c\) be the finite groupoid of TPK routes
with that profile and let \(A(c)=C^*(G_c)\). At matrix level \(s\), one takes the
cone
\[
 \mathcal C_s(c)=M_s(A(c))_+.
\]
A route \(r:c\to d\), represented by the partial isometry \(v_r\),
induces the CP map
\[
 T_r(a)=v_r^*av_r.
\]

\begin{theorem}[Functorial conic system of TPK paths]
\label{thm:cpe-s-system-tpk}
The assignments \(c\mapsto(\mathcal C_s(c))_{s\ge1}\) and
\(r\mapsto T_r\) respect identity, composition, matrix amplification, and
genealogical truncation. The CPM category generated by the finite fibers
is compact monoidal: the dual of a fiber is its conjugate space, and the
coevaluations and evaluations are the canonical cups and caps. The projective
limit of these systems preserves all compatible CP maps.
\end{theorem}

\begin{proof}
The complete positivity of \(a\mapsto v_r^*av_r\) is immediate from its
Kraus form. For composable routes,
\(v_{r_2r_1}=v_{r_2}v_{r_1}\), hence
\(T_{r_2r_1}=T_{r_1}T_{r_2}\) with the declared variance. The inclusions and
conditional expectations of the finite levels are UCP and form commuting
squares. The compactification of CPM in finite dimension is the standard construction
by conjugate spaces; truncation compatibility permits
passage to the limit coordinate by coordinate.
\end{proof}

Let \(\mathsf P_{\rm HMT}\) the set countable of codigos of routes
finite and let \(\mathsf C_{\rm HMT}\) the space of states with its contexto
complete. The evaluation
\[
 \operatorname{Ev}(\ulcorner r\urcorner,x)=T_r(x)
\]
preserves origin, destination, sheet, orientation, carry, boundary, and word of route.

\begin{corollary}[Genealogically faithful universality in the HMT domain]
\label{cor:cpe-universalidad-genealogica}
The evaluator \(\operatorname{Ev}\) is universal for the declared class of finite TPK transports and their compatible limits: every transport in that class has a code, composition is performed by concatenating codes, and reduction preserves the complete genealogy. This internal universality is not extended by decree to an untyped external class of targets.
\end{corollary}

\section{Identification of the Stinespring memory with the local record TPK}
\label{sec:cpe-stinespring-tpk}

In the nonadic expectation, the Stinespring dilation retains the discarded digit
in a basis \(\{e_a\}_{a=0}^8\). For a TPK state \(x\), let
\(\lambda(x,a)\) be the complete label of the continuation by the digit
\(a\). Define the isometry
\[
 W_x:e_a\longmapsto e_{\lambda(x,a)}
\]
from the memory minima of Stinespring toward the record local TPK\@. Let
\(F_xe_{\lambda(x,a)}=e_a\) the forgetful functor restringido to its imagen.

\Needspace{12\baselineskip}
\begin{theorem}[The TPK memory refines the minimal memory of Stinespring]
\label{thm:cpe-stinespring-tpk}
For every \(x\),
\[
 W_x^*W_x=I_9,
 \qquad F_xW_x=I_9.
\]
The family \((W_x)_x\) intertwines the inclusion of a new digit with the truncations. At depth \(m\), the tensor product of these isometries identifies the ancilla \((\mathbb C^9)^{\otimes m}\) with the TPK subregister that preserves the sequence of extensions; the additional sheet, carry, boundary, and route coordinates explain why TPK memory is a refinement rather than a second copy of the ancilla.
\end{theorem}

\begin{proof}
Continuations with distinct digits have orthogonal labels, so
\(W_x\) is isometric. The definition of \(F_x\) proves the second
identity. Composition of extensions concatenates the labels in the same
order as the tensor product of the ancillas; hence naturality under
truncation and the assertion at arbitrary depth.
\end{proof}
